\documentclass[11pt]{article}
\usepackage{amsmath,amssymb,amsthm,geometry,hyperref}
\geometry{margin=1in}
\newtheorem{theorem}{Theorem}
\newtheorem{lemma}{Lemma}
\newtheorem{corollary}{Corollary}
\newcommand{\C}{\mathbb C}
\newcommand{\Q}{\mathbb Q}
\newcommand{\len}{\operatorname{length}}
\newcommand{\ord}{\operatorname{ord}}
\title{An independent local-algebra verification of the Milnor-number
formula in manuscript family 059}
\author{Campaign 9000059: Engineer submission for independent review}
\date{November 5, 2026}
\begin{document}
\maketitle

\begin{abstract}
For the actual Weierstrass-prepared four-variable germs in the
multiplicities-four-and-five manuscript, we prove
\[
 \mu(f_{0,s})=\mu(f_{1,s})=24ab+6a+6b+51
\]
for every fixed complex \(s\ne0\) and integers \(a,b>1\) satisfying
\(ab>5(a+b)\).  The two explicit inequalities in its tube theorem
imply this condition.  The proof is a local intersection-length
calculation; it does not use a tube identification, a Seifert form,
or an unjustified replacement of a Weierstrass unit.  Exact rational
finite-jet calculations certify the three plane Jacobian-algebra
inputs, with a Nakayama argument proving that the truncations suffice.
This verifies a necessary numerical consequence, not the tube theorem
or either manuscript's ambient-homeomorphism assertion.  The analytic
proof is not kernel-formalized, and independent Reviewer certification
and novelty assessment are outstanding.
\end{abstract}

\section{Scope, sources, and the precise result}

Write M45 for the staged September 27, 2026 manuscript
\emph{Ambiently homeomorphic isolated hypersurface germs in
\(\C^4\) with multiplicities four and five}, and M23 for the staged
September 24, 2026 manuscript
\emph{Ambiently homeomorphic isolated hypersurfaces of multiplicities
two and three}.  They are independent arguments.

All 22 staged \texttt{.tex} files were read and inventoried.
\texttt{PROOF\_GRAPH.json} and \texttt{MATH\_STATE.json} contain
114 source-qualified assertions: all theorem, lemma, proposition,
corollary, and remark environments, the labeled displayed identities,
and the three central M45 construction definitions.  Source paths,
line intervals, original statements, and logical dependencies are
retained; merely citing another section is not treated as a proof
dependency.  None of these source assertions was certified merely
by being inventoried.

The local mission concerns M45, Section 03, and the numerical
consequence of \texttt{thm:tube} in Section 04.  Its equations are
\begin{align*}
 h(X,T)&=T(X-T^2)(X-2T^2),\\
 q_1(u,w)&=(u-w-w^8)(u-w-2w^8)(u-w^2),\\
 q_2(u,w)&=q_1(u,w)-u^2,\\
 \varphi(u,w)&=u^2+q_1(u,w),\\
 \psi(X,w)&=u(u-w^2)q_2(u,w),\qquad u=X+w .
\end{align*}
Let
\[
 F_0=X-T^2-q_1(T,w)=U_0R_0,\qquad
 F_1=T-\psi(X,w)=U_1R_1
\]
be the exact Weierstrass preparations in \(w\).  With
\[
 R_{i,s}(p,w)=s^4R_i(p,w/s),\quad
 C_{i,s}=\{yz=R_{i,s}(p,w)\},\quad
 f_{i,s}=h(p)+w(y^a+z^b),
\]
the functions are restricted to the indicated smooth ambient germs.
Here and throughout this report, \(\mu\) means the dimension of the
Jacobian quotient in the convergent local ring at the origin.

\begin{theorem}[Actual prepared germs]\label{thm:actual}
For every fixed \(s\in\C^*\) and every pair of integers \(a,b>1\)
with \(ab>5(a+b)\), both \(f_{i,s}\) have isolated critical point at
the origin and
\[
 \mu(f_{i,s})=24ab+6a+6b+51,\qquad i=0,1.
\]
This includes every pair satisfying the two strict inequalities in
M45's \texttt{thm:tube}, hence the eventual common parameter choice
in \texttt{lem:simultaneous-parameters}.
\end{theorem}

The assertion for every fixed \(s\ne0\) concerns germs: neighborhoods
may shrink with \(s\).  It asserts neither uniform isolation in the
manuscript's fixed outer tube nor any property at \(s=0\).
In particular, it does not establish \texttt{lem:local}.

\section{Local-algebra tools and the analytic input}

We use only local rings.  If \(g_1,\ldots,g_d\) have isolated common
zero in a regular local analytic ring of dimension \(d\), denote the
length of their quotient by \(I(g_1,\ldots,g_d)\).
Whenever the relevant intersections are finite,
\begin{equation}\label{eq:add}
 I(g_1,\ldots,g_{d-1},uv)
 =I(g_1,\ldots,g_{d-1},u)+I(g_1,\ldots,g_{d-1},v).
\end{equation}
Indeed the first \(d-1\) equations give a one-dimensional
Cohen--Macaulay ring \(A\), in which \(u,v\) are nonzerodivisors.
The exact sequence
\[
 0\longrightarrow A/(u)\xrightarrow{\ \cdot v\ }A/(uv)
       \longrightarrow A/(v)\longrightarrow0
\]
proves the formula.  Finiteness of the two intersections on the
right also implies finiteness on the left, so the formula is not
being used to assume isolation.  The regular-sequence and
Cohen--Macaulay facts used here are precisely
\cite[Tags 00NQ and 02JN]{stacks}.

\begin{lemma}[A monic free quotient]\label{lem:free}
For any commutative coefficient ring \(B\), any \(\rho\in B\), and
integers \(a,b>1\), the polynomial quotient
\[
 B[y,z]/(yz-\rho,\ y^a+z^b)
\]
is free over \(B\), with basis
\[
 1,y,\ldots,y^{a-1},z,\ldots,z^b.
\]
If \(B\) is local Artinian and \(\rho\) is in its maximal ideal,
this is also the corresponding local power-series quotient.
\end{lemma}
\begin{proof}
For the weighted monomial order of weights \((b,a)\), with \(y\)
first in ties, a monic basis is
\[
 g_1=yz-\rho,\quad g_2=y^a+z^b,\quad
 g_3=z^{b+1}+\rho y^{a-1}.
\]
Its leading monomials are \(yz,y^a,z^{b+1}\).
The three pair reductions are
\[
 y^{a-1}g_1-zg_2=-g_3,\qquad
 z^bg_1-yg_3=-\rho g_2,
\]
\[
 z^{b+1}g_2-y^ag_3=z^bg_3-\rho y^{a-1}g_2.
\]
Thus monic division gives exactly the displayed basis, without any
division by a coefficient of \(B\).
In the Artinian case,
\(y^{a+b}=-\rho^b\) and \(z^{a+b}=-\rho^a\), so both variables
are nilpotent.  Localization and completion therefore add no points
and change no length.
\end{proof}

The essential analytic step is the following small-perturbation
calculation.  It applies to analytic coefficients, not merely to
polynomial approximations of the prepared equation.

\begin{lemma}[Exact polar intersection]\label{lem:polar}
Let \(p=(v,u)\), let \(h(p)\) have an isolated critical point, and
suppose for every nonconstant analytic \(p\)-arc
\[
 \nu(\nabla h)\le L\nu(p),\qquad
 \nu(p)=\min(\nu v,\nu u).
\]
Let \(R(p,w)=w^k+\sum_{j<k}A_j(p)w^j\), with \(A_j(0)=0\).
Put \(m=a+b,\ n=ab,\ N=n-m\), and assume \(N/m>L\).
For any \(\xi\ne0\), the four equations
\begin{align}
 R-\xi t^m&=0,\qquad mR+nwR_w=0,\label{eq:polar-geometry}\\
 \nabla h+\lambda(b/\xi)wt^N\nabla_pR&=0\label{eq:polar-gradient}
\end{align}
in \((p,w,t)\) have isolated intersection at the origin of
length \(km\mu(h)\), for every finite complex \(\lambda\).
\end{lemma}
\begin{proof}
The equation \(mR+nwR_w=0\) is a degree-\(k\) Weierstrass
equation in \(w\): its leading coefficient is \(m+nk\ne0\),
and all lower coefficients vanish at \(p=0\).
On an arc with \(p\not\equiv0\), it forces
\(\nu w\ge\nu p/k\), since otherwise its leading term has
strictly smaller order than every other term.  The first equation
in \eqref{eq:polar-geometry} then implies
\(m\nu t=\nu R\ge\nu p\).  Consequently the perturbing vector
in \eqref{eq:polar-gradient} has order at least
\[
 \nu w+N\nu t\ \ge\ (1/k+N/m)\nu p\ >\ L\nu p .
\]
It cannot cancel the lowest-order component of \(\nabla h\).
The same reasoning applies to arcs with variable \(\lambda\)
near any finite \(\lambda_0\), because \(\lambda\) has no pole.
If \(p\equiv0\), the second equation gives \(w\equiv0\), and
the first then gives \(t\equiv0\).

For completeness, this arc argument gives uniform local isolation
in \(\lambda\), not just an assertion about individual fibers.
An analytic component off the parameter axis would, after a generic
linear slice, contain a one-dimensional analytic germ off that
axis; a Puiseux parametrization supplies the forbidden arc.
These are the local-parametrization and Puiseux statements
in \cite[Theorems 6.7.4 and 6.7.6]{lebl}.

The analytic quotient defined by the four equations near
\((0,\lambda_0)\) is therefore finite over
\(\C\{\lambda-\lambda_0\}\) and has support the parameter axis.
It is a dimension-one complete intersection in a regular local
ring of dimension five.  It is Cohen--Macaulay, and the base is a
regular local ring of dimension one with zero-dimensional fiber.
Thus the precise hypotheses of miracle flatness
\cite[Tag 00R4]{stacks} hold.  The finite flat module is free locally,
so its fiber length is constant.  Connecting \(0\) to any finite
\(\lambda\) proves equality of the lengths.

At \(\lambda=0\), first quotient by \(\nabla h\), giving an Artinian
ring of length \(\mu(h)\).  The second equation in
\eqref{eq:polar-geometry}, divided by \(m+nk\), is monic of degree
\(k\) in \(w\); the first is monic of degree \(m\) in \(t\), up
to its nonzero constant coefficient.  Successive Weierstrass
division gives free ranks \(k\) and \(m\).  The length is
\(km\mu(h)\).
\end{proof}

\section{The general intersection-length formula}

\begin{theorem}\label{thm:general}
In the situation of Lemma~\ref{lem:polar}, assume in addition that
\(R_v(0)\ne0\), that \(d=h|_{\{R=0\}}\), in coordinates \((u,w)\),
has an isolated critical point, and that \(d(u,0)\not\equiv0\).
Then on the smooth germ \(yz=R(p,w)\),
\[
 f=h(p)+w(y^a+z^b)
\]
has an isolated critical point and
\[
 \mu(f)=\mu(d)+(kab+a+b)\mu(h).
\]
\end{theorem}
\begin{proof}
Solve \(R(v,u,w)=q\) exactly for \(v\), and put
\(H(u,w,q)=h(v(u,w,q),u)\).  No unit is suppressed.
All derivatives of \(R\) below are evaluated at this exact solution.
Implicit differentiation gives
\begin{equation}\label{eq:chain}
 H_q=h_v/R_v,\qquad
 H_u=h_u-(R_u/R_v)h_v,\qquad
 H_w=-R_wH_q.
\end{equation}
The function in the four smooth coordinates is
\[
 f=H(u,w,yz)+w(y^a+z^b).
\]
Set
\[
 A=ay^a-bz^b,\quad
 F_1=H_u,\quad F_2=H_w+y^a+z^b,\quad
 F_3=zH_q+aw y^{a-1}.
\]
Then \(F_1=f_u,F_2=f_w,F_3=f_y\), and
\(yF_3-zf_z=wA\).  Define
\[
 L_A=I(F_1,F_2,F_3,A),\quad
 L_w=I(F_1,F_2,F_3,w),\quad
 L_z=I(F_1,F_2,F_3,z).
\]
We first compute these finite lengths, without assuming isolation
of \(f\).

Write \(\ell=\ord d(u,0)-1\).  At \(z=0\), the three remaining
equations are \(d_u,d_w+y^a,aw y^{a-1}\), so
\begin{equation}\label{eq:Lz}
 L_z=a\ell+(a-1)\mu(d).
\end{equation}
Indeed the \(w=0\) factor has length \(a\ell\), by monic
division in \(y\), and the \(y=0\) factor has length \(\mu(d)\).

At \(w=0\), the last equation is \(zH_q\).
Its \(z=0\) contribution again has length \(a\ell\).
On \(H_q=H_u=0\), equations \eqref{eq:chain} are exactly
\(\nabla h=0\); also \(H_w=0\).  The remaining equations are
\(yz=R(p,0)\) and \(y^a+z^b=0\).
Apply Lemma~\ref{lem:free} over
\(B=\C\{p\}/(\nabla h)\), with \(\rho=R(p,0)\) nilpotent.
This proves
\begin{equation}\label{eq:Lw}
 L_w=a\ell+(a+b)\mu(h).
\end{equation}

To compute \(L_A\), make the finite free cover
\[
 y=Y^b,\qquad z=Z^a.
\]
Its degree is \(ab=n\): the monomials \(Y^iZ^j\),
\(0\le i<b,\ 0\le j<a\), form a free basis.
The pulled-back equation \(A\) is
\(aY^n-bZ^n\), a product of \(n\) distinct linear factors.
On a factor \(Y=\eta Z\), where \(\eta^n=b/a\), set \(t=Z\)
and \(\xi=\eta^b\).  Thus
\[
 y=\xi t^b,\quad z=t^a,\quad
 q=\xi t^m,\quad \xi^a=b/a.
\]
All occurrences of \(H\) in the following formulas are at
\((u,w,\xi t^m)\).  The three equations become
\[
 H_u,\qquad H_w+(m/a)t^n,\qquad
 t^a\bigl(H_q+(b/\xi)wt^N\bigr).
\]
The factor \(t^a\) contributes \(a\mu(d)\).
Put \(D=H_q+(b/\xi)wt^N\).  Modulo \(D\), the middle equation
can be replaced, using \eqref{eq:chain}, by
\[
 t^N\bigl((m/a)t^m+(b/\xi)wR_w\bigr).
\]
The \(t^N\) factor contributes \(Nk\mu(h)\):
at \(t=0\), the equations \(H_u=H_q=0\) are the gradient
equations of \(h\) on \(R=0\), and monic division in \(w\)
gives rank \(k\) over its Jacobian algebra.

For the last factor, use \(R=\xi t^m\) and multiply by the unit
\(a\xi\); its equation is \(mR+nwR_w=0\).
Using \eqref{eq:chain}, the other two equations are equivalent to
\[
 \nabla h+(b/\xi)wt^N\nabla_pR=0.
\]
Lemma~\ref{lem:polar} at \(\lambda=1\) gives length \(km\mu(h)\).
The length on each of the \(n\) linear factors is therefore
\[
 a\mu(d)+Nk\mu(h)+mk\mu(h)=a\mu(d)+nk\mu(h).
\]
All these lengths are finite.  Repeated application of
\eqref{eq:add}, followed by division by the free-cover degree
\(n\), proves
\begin{equation}\label{eq:LA}
 L_A=a\mu(d)+kab\mu(h).
\end{equation}

Now \(I(F_1,F_2,F_3,wA)=L_w+L_A\) is finite.  Since its ideal
also equals \((F_1,F_2,F_3,zf_z)\), the Jacobian ideal is
zero-dimensional and \(f\) is isolated.  Applying
\eqref{eq:add} one last time gives
\[
 \mu(f)+L_z=L_w+L_A.
\]
Substituting \eqref{eq:Lz}--\eqref{eq:LA} proves the assertion.
In particular, the terms involving the slice order \(\ell\)
cancel exactly.
\end{proof}

\section{Certified inputs for the actual germs}

The exact polynomial graph pullbacks are
\[
 d_0(u,w)=u q_1(u,w)q_2(u,w),\qquad
 d_1(X,w)=\psi(X,w)(X-\psi(X,w)^2)(X-2\psi(X,w)^2).
\]
The verifier constructs these expressions from their factors,
using rational coefficients only.

\subsection*{Why the finite jets prove the local dimensions}
For a polynomial \(g(x,w)\), the verifier computes
\[
 L_D(g)=\dim_\Q \Q[x,w]/(g_x,g_w,(x,w)^D)
\]
by exact rational row reduction.  Columns are all monomials of
total degree below \(D\); rows are all monomial multiples of the
two derivatives that survive modulo \((x,w)^D\).
This is the entire ideal in that finite quotient, not a sampled
set of relations.

If \(L_D=L_{D+1}\), put \(M=\C\{x,w\}/J(g)\).  The natural
surjection \(M/\mathfrak m^{D+1}M\to M/\mathfrak m^DM\)
then has zero-dimensional kernel, hence
\(\mathfrak m^DM=\mathfrak m^{D+1}M\).
The module \(\mathfrak m^DM\) is finitely generated over the
Noetherian local ring, even before \(M\) is known Artinian.
Nakayama's lemma \cite[Tag 00DV, part (2)]{stacks} gives
\(\mathfrak m^DM=0\).  Thus \(L_D\) is the exact local length
and proves isolation.  Extending rational scalars to \(\C\)
does not change these finite ranks.

The certificate records the following exact outputs:
\[
\begin{array}{c|c|c|c}
g&D&D(D+1)/2-L_D&L_D=L_{D+1}\\ \hline
h&6&15&6\\
d_0&24&249&51\\
d_1&24&249&51
\end{array}
\]
The next-cutoff ranks are respectively \(22,274,274\).
For the plane pullbacks, function terms of degree at least 26
cannot affect either tested derivative quotient.  Equality of the
two quotient lengths supplies the missing justification that an
arbitrary truncation would lack.
Both \(d_i(u,0)\) have initial term \(-u^6\), so the slice order
is six.  These computations independently recover the numerical
part of M45's \texttt{lem:plane-data}; they do not verify its
Seifert-form or monodromy assertions.

As a separate bounded operation, the approved Jacobian adapter
computed the exact lexicographic basis
\[
 J(h)=\left(X^2-\tfrac72T^4,\ XT-\tfrac32T^3,\ T^5\right).
\]
This ideal is supported only at the origin.  Its six standard
monomials are \(1,T,T^2,T^3,T^4,X\).
The dedicated JSON certificate preserves the operation ID,
exact payload, full typed response, server version, protocol, and
successful adapter status.  The standalone verifier replays the
basis and the local jet calculation.  No adapter result concerning
a global ideal is used as a substitute for either four-variable
local algebra.

\subsection*{Analytic units and all nonzero \(s\)}
Direct polynomial identities give
\[
 F_0(0,0,w)=w^4(1+w^7)(1+2w^7),
\]
\[
 F_1(0,0,w)=w^4\bigl((1-w)-2w^{15}(1-w)^2\bigr).
\]
The two factors after \(w^4\) have constant term one, and
\(F_{0,X}(0)=F_{1,T}(0)=1\).
The hypotheses of Weierstrass preparation
\cite[Theorem 6.2.3]{lebl} therefore hold exactly:
\[
 R_i(0,w)=w^4,\qquad R_{0,X}(0)=R_{1,T}(0)=1.
\]
For each fixed \(s\ne0\), \(R_{i,s}\) is still monic of degree
four with lower coefficients vanishing at \(p=0\), and the indicated
base derivative is \(s^4\ne0\).
Its root surface is the exact graph with \(w\) replaced by \(w/s\).
Consequently its plane pullback is analytically equivalent to \(d_i\)
by an invertible linear change in \(w\), and its Jacobian length is
51.  The proof of Theorem~\ref{thm:general} applies directly to
\(R_{i,s}\); it never sets \(U_i\) equal to one, nor truncates
these prepared analytic functions.

\subsection*{Gradient and admissible parameters}
The derivatives of \(h\) are
\[
 h_X=T(2X-3T^2),\qquad h_T=X^2-9XT^2+10T^4.
\]
If \(\nu X<2\nu T\), the initial term of \(h_T\) is \(X^2\);
if \(\nu X>2\nu T\), the initial term of \(h_X\) is \(-3T^3\).
Both give \(\nu\nabla h\le4\min(\nu X,\nu T)\).
In the equality case, writing \(X=cT^2+\cdots\), the identity
\[
 (2c-15)(2c-3)-4(c^2-9c+10)=5
\]
excludes simultaneous cancellation.  Arcs on either coordinate
axis satisfy the bound directly.  Hence \(L=4\).

Let \(\alpha=4b/(a+b)\) and \(\beta=4a/(a+b)\).
The exact identity
\[
 \frac{ab-a-b}{a+b}-4
 =\frac12\left(\frac{(a-2)\alpha}{4}-4\right)
  +\frac12\left(\frac{(b-2)\beta}{4}-4\right)
\]
shows that both strict tube inequalities imply \(N/m>4\).
The general formula, with \(k=4,\mu(h)=6,\mu(d_i)=51\),
now proves Theorem~\ref{thm:actual}.  The universal inequality is
checked symbolically, not inferred from finitely many parameter
examples.  No claim about the manuscript's unspecified
``sufficiently small'' threshold is needed: the local conclusion
holds for every fixed nonzero \(s\).

\section{What is, and is not, certified by the computation}

\texttt{verify.py} checks the source inventory, the exact plane
Jacobian lengths with a termination certificate, the graph-unit
identities and base derivatives, the gradient noncancellation
identity, the monic-quotient reductions, the parameter inequality,
and the polar and final length identities.  It is rerunnable using
Python with SymPy; the plane rank computation itself uses only the
standard-library \texttt{Fraction} class.  It needs neither a
Jacobian server nor Singular when replaying the stored result.
The software-side distinction between a local and a global
Jacobian quotient is explicitly documented in the primary
Singular manual \cite[Section A.4.1]{singular}.

The analytic arc argument, finite-flatness application, and use of
intersection-length additivity are the mathematical proof above,
not outcomes of a proof kernel.  The independent Reviewer must
check these arguments and the fidelity of the cited foundational
statements.  Citations have been read by the Engineer and recorded
as assumptions; this is not an independent citation attestation.

Compared with the prior state, which contained no verified local
Milnor computation, this supplies an all-parameter local-algebra
theorem rather than matching examples or a finite exclusion.
A focused literature search did not locate this exact formula;
that does not establish novelty.  Attempts to obtain the original
Milnor monograph from two public endpoints did not yield a readable
source, so no unread theorem locator from that book is relied upon.
Here the Milnor number is the local-algebra invariant itself.

The equality of Milnor numbers neither proves integral Seifert
congruence nor identifies the custom large tube with a local Milnor
representative.  Acyclicity alone would not justify that
identification.  All those obligations, the existential M45 main
theorem, and the independent M23 manuscript remain open in the
campaign.  No refutation of either existential assertion is claimed.

\begin{thebibliography}{9}
\bibitem{stacks}
The Stacks Project Authors, \emph{The Stacks Project}, read November
5, 2026:
\href{https://stacks.math.columbia.edu/tag/00DV}{Tag 00DV},
Lemma 10.20.1(2) (Nakayama);
\href{https://stacks.math.columbia.edu/tag/00NQ}{Tag 00NQ},
Lemma 10.106.3 (regular local rings are Cohen--Macaulay);
\href{https://stacks.math.columbia.edu/tag/02JN}{Tag 02JN},
Lemma 10.104.2 (regular sequences and dimension);
\href{https://stacks.math.columbia.edu/tag/00R4}{Tag 00R4},
Lemma 10.128.1 (miracle flatness).
\bibitem{lebl}
J.~Lebl, \emph{Tasty Bits of Several Complex Variables},
version 4.4, May 31, 2026, ISBN 979-8866906680.
Theorems 6.2.3 and 6.2.5 (preparation and division);
Theorems 6.7.4 and 6.7.6 (local parametrization and Puiseux).
\url{https://www.jirka.org/scv/scv.pdf}.
\bibitem{singular}
The Singular authors, \emph{Singular 4.4 Manual},
Section A.4.1, ``Milnor and Tjurina number'',
\url{https://www.singular.uni-kl.de/Manual/4-4/sing_960.htm}.
See also the local/global contrast in Section 2.3.4,
\url{https://www.singular.uni-kl.de/Manual/latest/sing_9.htm}.
\end{thebibliography}
\end{document}
